%=============================================================================!
% 14.5  WHAT A BAROSTAT MUST DO                                               !
%=============================================================================!
\section{What a barostat must do}
\label{sec:pr-npt}

A barostat lets the volume of the box change so that the run samples the
ensemble at constant pressure. This section sets out what that ensemble
requires, finds the exact distribution of the volume for an ideal gas and
its spread for a liquid, which are the tests every barostat of the
chapter must pass, and fits a barostat into the loop of Chapter~13.

\subsection*{The ensemble at constant pressure}

A system that exchanges volume with a bath at the pressure $P_0$ has the
weight $\mathrm{e}^{-\beta(H + P_0V)}$ (\cref{sec:sm-other}). Grouping the
states by their volume gives the probability of each volume,
\begin{equation*}
   \mathcal{P}(V) \propto \mathrm{e}^{-\beta P_0V}\,Z(N, V, T) ,
\end{equation*}
the canonical partition function of the box at that volume weighted by
$\mathrm{e}^{-\beta P_0V}$. A barostat must therefore do two things that a
thermostat does not: move the volume, and move it so that it takes each
value with this probability. Its mean pressure should be $P_0$, and the
spread of its volume should be the one this distribution has.

\subsection*{The ideal gas exactly}

For atoms that do not interact, $U = 0$ and \cref{sec:pr-periodic} gives
$Z = C\,V^N$, so
\begin{equation*}
   \mathcal{P}(V) \propto V^N\mathrm{e}^{-bV} , \qquad b = \frac{P_0}{\kB T} .
\end{equation*}
The integrals $\int_0^\infty V^n\mathrm{e}^{-bV}\dd V = n!/b^{n+1}$, found by
integrating by parts $n$ times (\cref{ex:pr-gamma}), give the mean and the
variance, each the integral with one or two more powers of $V$ divided by
the integral with none:
\begin{align*}
   \ens{V} &= \frac{(N + 1)!/b^{N+2}}{N!/b^{N+1}} = \frac{N + 1}{b}
   = \frac{(N + 1)\kB T}{P_0} ,\\
   \ens{V^2} &= \frac{(N + 2)!/b^{N+3}}{N!/b^{N+1}}
   = \frac{(N + 2)(N + 1)}{b^2} , \qquad
   \ens{V^2} - \ens{V}^2 = \frac{N + 1}{b^2} .
\end{align*}
The mean is the ideal gas law's with $N + 1$ in place of $N$, a difference
that vanishes for large $N$, and the relative spread is $1/\sqrt{N + 1}$.
For 20 atoms at \SI{135}{\kelvin} and $P_0 = \SI{1e-4}{\electronvolt\per
\angstrom\cubed}$, the mean volume is \SI{2443}{\angstrom\cubed} and its
spread \SI{533}{\angstrom\cubed}: a small gas whose volume fluctuates by a
fifth, a severe test (\texttt{virial.ideal\_gas\_volume\_density} gives
the whole density).

\subsection*{A liquid's volume}

Chapter~12 found that the volume at constant pressure fluctuates with the
variance $-\kB T\,\dd\ens{V}/\dd P$ (\cref{ex:sm-volume}). With the
compressibility of \cref{sec:pr-periodic}, $\dd\ens{V}/\dd P =
-\ens{V}\kappa_T$, so
\begin{equation}
   \ens{V^2} - \ens{V}^2 = \kB T\,\ens{V}\,\kappa_T .
   \label{eq:pr-volume}
\end{equation}
A soft material swings its volume more than a stiff one. To test a
barostat on the liquid this needs $\kappa_T$ found without one. Three
runs at fixed volume around the volume the liquid takes at
\SI{0.1}{\giga\pascal}, three starts each of \SI{100}{\pico\second} under
CSVR at \SI{135}{\kelvin}, give
\begin{center}
\small
\begin{tabular}{@{}lccc@{}}
\toprule
$V$ / \si{\angstrom\cubed} & 12\,112 & 12\,301 & 12\,477\\
$\ens{P}$ / \si{\giga\pascal} & $0.1109 \pm 0.0004$ & $0.1000 \pm 0.0005$
   & $0.0904 \pm 0.0013$\\
\bottomrule
\end{tabular}
\end{center}
with the standard error of the three starts. The parabola through the
three points has, at \SI{12306}{\angstrom\cubed}, the mean volume of the
runs at constant pressure of \cref{sec:pr-rescaling,sec:pr-piston}, the
slope $\dd P/\dd V$ that gives $\kappa_T = -1/(V\,\dd P/\dd V) = (1.46 \pm
0.11)$\,\si{\per\giga\pascal}, the error found by redrawing the three
pressures within their standard errors; the liquid's bulk modulus is a
quarter of the crystal's. \Cref{eq:pr-volume} then predicts a spread of
\SI{183}{\angstrom\cubed}, 1.5\% of the volume.

\subsection*{The barostat in the loop}

\texttt{barostats.run} is Chapter~13's loop with one addition. At the
start of each step, after the thermostat has acted, the barostat reads
the instantaneous pressure tensor, \cref{eq:pr-tensor}, from the
velocities and the virial the last force call left, and changes the cell,
scaling the positions with it so that their fractional coordinates are
kept, and, for some barostats, the velocities; the forces are then
computed again in the new cell and the step continues as velocity
Verlet. The neighbour list is kept across such changes while no pair can
have been missed. If the strain since the list was built stretches or
shrinks no separation by more than the fraction $e$ of its length (for a
box scaled by the factor $1 + e$ along every edge, $e$ itself), and no atom
has moved further than $u$ from where the strain alone would have carried
it, a pair now within $r_{\mathrm c}$ was within $r_{\mathrm c} + 2u$ after
the strain alone, and so within $(r_{\mathrm c} + 2u)/(1 - e)$ when the
list was built. The list of \cref{sec:md-neighbours} therefore holds while
$(r_{\mathrm c} + 2u)/(1 - e) \le r_{\mathrm c} + r_{\mathrm{skin}}$; with the
box fixed, $e = 0$ and this is the rule $u \le \frac12r_{\mathrm{skin}}$ of
that section. The run records the work, the change of
$K + U + P_0V$ that the barostat makes, beside the heat of the
thermostat; the effective energy $K + U + P_0V$ less the heat and the work
changes only through the integration error, as \cref{eq:th-effective} did
for a thermostat.

\begin{keyidea}
A barostat must sample $\mathcal{P}(V) \propto \mathrm{e}^{-\beta
P_0V}Z(N, V, T)$: the mean pressure $P_0$, and the right spread of volume.
For an ideal gas $\mathcal{P}(V) \propto V^N\mathrm{e}^{-\beta P_0V}$, with
mean $(N + 1)\kB T/P_0$ and relative spread $1/\sqrt{N + 1}$; for a liquid
the variance of $V$ is $\kB T\ens{V}\kappa_T$, with $\kappa_T$ measured at
fixed volumes.
\end{keyidea}

\begin{selfcheck}
By what fraction does the volume of 1000 ideal-gas atoms fluctuate at
constant pressure?
\end{selfcheck}
